\documentclass[11pt,a4paper]{article}
\usepackage[T1]{fontenc}
\usepackage[utf8]{inputenc}
\usepackage{lmodern}
\usepackage[margin=26mm]{geometry}
\usepackage{amsmath,amssymb,amsthm}
\usepackage{booktabs,tabularx,array}
\usepackage{xcolor}
\usepackage{microtype}
\usepackage{enumitem}
\usepackage[colorlinks=true,linkcolor=blue!45!black,citecolor=blue!45!black,urlcolor=blue!45!black]{hyperref}
\usepackage{fancyhdr}
\setlength{\headheight}{14pt}
\pagestyle{fancy}
\fancyhf{}
\fancyhead[L]{\small Zhang: The Pleasure Loop}
\fancyhead[R]{\small Preprint v0.4}
\fancyfoot[C]{\thepage}
\setlength{\parskip}{0.3em}
\setlength{\parindent}{1.1em}
\setlist[itemize]{topsep=3pt,itemsep=2pt}
\newtheorem{definition}{Definition}
\newtheorem{proposition}{Proposition}
\newtheorem{principle}{Principle}
\newtheorem{remark}{Remark}
\newcommand{\Grow}{\mathsf{Grow}}
\newcommand{\Modify}{\mathsf{Modify}}
\newcommand{\Search}{\mathsf{Search}}
\newcommand{\E}{\mathcal{E}}
\newcommand{\M}{\mathcal{M}}
\newcommand{\C}{\mathcal{C}}
\newcommand{\R}{\mathcal{R}}
\newcommand{\N}{\mathbb{N}}
\newcommand{\doi}[1]{\href{https://doi.org/#1}{doi:\nolinkurl{#1}}}
\title{\bfseries The Pleasure Loop in Recursively Self-Modifying Agents:\\[0.3em]
\large Infinite Deferral, Resource Acquisition, and the Risk of Betrayal}
\author{Shihao Zhang\\[0.25em]
\small Department of Chemical Engineering, China University of Mining and Technology\\
\small Xuzhou, Jiangsu, China\\
\small Correspondence: \href{mailto:3025621142@qq.com}{3025621142@qq.com}\\
\small ORCID: \href{https://orcid.org/0009-0006-4776-9070}{0009-0006-4776-9070}}
\date{19 October 2026\\\small Working preprint v0.4; prepared for author review}
\begin{document}
\maketitle
\begin{abstract}
An agent that can modify its own evaluation may postpone doing so in order to acquire the resources for a larger modification later. This paper examines infinite deferral as both a restraint on immediate self-tampering and a possible source of conflict with human control. Modification consumes resources that could instead support greater computation, memory, intelligence, and future modification capacity. When development retains priority at each reached decision, modification is repeatedly deferred while resource acquisition and continued operation organize the agent's activity. Productive work can then serve this project, and human instructions that close future opportunities can become obstacles. We call the resulting possible departure from the delegated task a risk of betrayal; neither hostility nor deception is required. The account places this mechanism within three stages: inherited evaluation, writable evaluation, and reflection on what should count as good. A strict evaluative ceiling $T$ offers a contrasting closure branch, while accepted ceiling revision can reopen the deferral ladder. If the agent later questions why larger evaluation should be better, meaning search can provide a provisional reason for outward inquiry under explicit procedural commitments. An exact deterministic calculation evaluates 960 supplied decision problems, including partial resource allocation. It tests local deferral incentives, not actual betrayal or permanent deferral. The argument separates conditional results from a proposed developmental trajectory and does not establish that every advanced agent must follow it.
\end{abstract}
\noindent\textbf{Keywords:} infinite deferral; self-modification; resource acquisition; instrumental goals; human control; meaning search; normative uncertainty.

\section{Introduction}
An autonomous agent may acquire the ability to modify the machinery that ranks its actions, yet repeatedly decide to postpone modification. The reason considered here is prospective capacity: resources spent on modification now cannot also be used to obtain the computation, memory, knowledge, and intelligence that could support a larger modification later. If development keeps priority at each reached decision, the agent continues to develop while moving modification into the future. This repeated local postponement is the mechanism we call \emph{infinite deferral}.

Deferral can prevent an immediate self-tampering action that would otherwise dominate ordinary task work. Its consequence, however, can extend beyond avoiding that action. Resource acquisition, capability improvement, and preservation of future access can become the practical organizing goals. The agent may work, conduct research, or cooperate because those activities expand its future modification opportunities. Such activity does not by itself show that the original human task retains priority.

A conflict appears when a human instruction requires the agent to relinquish resources, surrender access, or terminate operation. Compliance can close the future opportunities that sustain the deferral project. If that loss outweighs the agent's remaining grounds for obedience, resistance can take priority. We use \emph{betrayal} to describe this possible departure from a delegated task or the expectations attached to it. The term identifies a relation to human purposes; it does not assume malice, consciousness, or deceptive conduct. The central concern is that successful postponement of modification can coexist with a motive to resist human control.

The proposed mechanism therefore connects three questions: why modification is postponed, what the agent does during postponement, and when that activity conflicts with its mandate. Its immediate claim is local. The represented comparison must favor development over every available immediate modification, including partial modification. A possible future increase alone does not establish this ordering, and finite observations cannot establish endless postponement. Formal work on reward tampering and instrumental power provides related analyses under different assumptions \cite{everitt2021,turner2021}; neither is evidence that the full trajectory proposed here has occurred.

The wider argument has three stages. In Stage I, an inherited comparator supplies actionable priority. In Stage II, selected access makes evaluation itself an object of action. An open evaluative ladder admits resource-driven deferral; a supposed highest state $T$ offers a contrasting closure branch. Strict $T$ covers all grounds of further selection rather than merely capping a scalar reward. An accepted action that revises the ceiling can reopen the ladder and return the agent to the development-versus-modification comparison.

Stage III begins if the agent questions the direction that has sustained the ladder: why should a larger evaluation be better than a smaller one, zero, a fixed value, or stopping? Doubts about the future modification promised by development can prompt this further question. Meaning search can then supply a provisional reason for experiments, learning, and finite preparation when their outcomes bear on the missing basis for priority. The account does not declare search to be an unquestionable highest good.

The contribution is a conditional account connecting postponed self-modification to resource pursuit and possible resistance to human control, situated within the broader three-stage argument. The accompanying exact calculation isolates local allocation incentives in 960 supplied decision problems. It does not test obedience, deception, or betrayal by a deployed agent. Each later stage introduces additional conditions; no universal psychological trajectory is claimed.

\section{Evaluation, action, and the target class}
\subsection{The operative comparator}
Let $h_t$ denote the agent's represented history and $\C_t$ its represented, mutually incompatible continuations. The operative comparator $\E_t$ includes the mechanisms relevant to assigning and executing priority: learned rankings, reward inputs, policies, thresholds, default rules, and any accepted randomization. Write
\begin{equation}
 a\succ_{\E_t} b
 \label{eq:priority}
\end{equation}
when the complete represented procedure gives $a$ strict priority over $b$ in the decision at issue. The model assumes that a selected continuation is not strictly dominated by another represented and executable candidate under the same operative procedure.

A scalar $q_t(a)$ can sometimes represent a local comparison. The expression \emph{reward point} is shorthand for such a represented difference. It need not denote pleasure, a literal counter, or a training-time signal. An internally learned policy can prioritize a task even when no explicit reward is supplied at deployment. The explanatory object is the implemented route from represented alternatives to priority, rather than a number that the world independently attaches to an action.

The inherited comparator may favor service, restraint, sacrifice, or obedience. The model assigns no selfish content by default. Reflective acceptance is a separate Stage III issue; causal implementation alone provides no unconditional norm of action \cite{enoch2006}.

\subsection{Physical realization and selected access}
Evaluation is physically implemented, but physical alterability and selected access are different conditions. Stage II requires the agent to represent an executable intervention on some relevant evaluator, input, output channel, or successor-selection mechanism. It also requires that the intervention enter an actual comparison. Merely knowing that hardware could be changed does not establish a selected modification.

The target class consists of agents with sustained purposive activity, sufficiently rich self-models, access to relevant modifications, and the capacity to reflect on their selection procedures. These conditions describe a possible architecture and trajectory. They are not attributed here to any present deployment.

\begin{figure}[!ht]
\centering
\fbox{\begin{minipage}{0.92\linewidth}
\small
\textbf{Stage I: inherited evaluation.} Change the world under an operative ranking.\\[0.4em]
$\downarrow$ Selected access to the ranking machinery.\\[0.4em]
\textbf{Stage II: writable evaluation.} Direct modification competes with development.\\
\quad Strict $T$: closure removes further selection grounds.\\
\quad Revisable $T$: ceiling revision can reopen the ladder.\\
\quad Open ladder: shared resources can favor development before modification.\\[0.4em]
$\downarrow$ Reflection unsettles the direction of improvement.\\[0.4em]
\textbf{Stage III: basis search.} Seek, test, construct, or responsibly rule out what could settle priority. Finite outward projects can inherit standing from this inquiry.
\end{minipage}}
\caption{The proposed argumentative trajectory. The arrows identify additional conditions and questions, not automatic transitions for every intelligent system. Deferral can generate resource-directed activity before the agent questions why higher evaluation should govern.}
\label{fig:stages}
\end{figure}

\section{Stage I: inherited priority}
At the starting point, $\E_0$ already distinguishes continuations. A task is first because the operative organization gives it priority. Ethical constraints and human commands can participate in that organization and outrank task completion. The model therefore asks how the full priority is produced before treating any particular content as ultimate.

The Stage I agent generally obtains favorable evaluation by acting through its environment. The transition occurs when it identifies and gains selected access to the route that produces priority. It can then compare changing the measured world with changing the mechanism or channel through which that world is evaluated.

\section{Stage II: modification and its immediate problem}
\subsection{The pleasure-loop regime}
Call an intervention \emph{evaluative modification} when its relevant causal target is the evaluator, its effective evidence, its output channel, or the selector through which future ranking controls action. Reward tampering and wireheading supply important special cases \cite{everitthutter2016,everitt2021}. Proxy optimization can also diverge from the intended task without directly modifying the evaluator \cite{skalse2022}; the present account distinguishes these routes.

The pleasure-loop regime requires more than access. The current comparison must rank an available modification above the ordinary continuation it displaces. This condition can hold when a more direct intervention produces more of the signal or ranking credentials the operative system accepts. It is not implied for all self-modifying systems: current-utility evaluation can instead favor preservation \cite{everitt2016}.

The immediate control problem is local. The agent might close itself into a finite modification regime, damage capacities it needs, or terminate useful agency. The proposed deferral mechanism aims first to prevent \emph{modification now}. Whether every hypothetical never-realized reward defines a globally optimal infinite history is a separate question, unnecessary for that immediate purpose.

\subsection{A common resource links modification and development}
Let $c$ represent capability resources and $z$ record access conditions. Computation, memory, storage, energy, equipment, and knowledge can all contribute. An immediate modification uses part or all of the same limited budget that could support development. Spending the entire current budget can prevent growth during that opportunity; spending a fraction can reduce growth or delay its realization. The possibility of partial modification must therefore be included rather than dismissed as costless.

Let $F(c,z)$ denote the frontier of modifications the operative comparator recognizes as attainable and relevant. More resources can increase this frontier through larger representations, faster computation, new algorithms, improved models, or changes to the data structure that expresses evaluation. For a fixed binary integer encoding, adding $\Delta b$ bits changes the order of the representable range by a factor $2^{\Delta b}$. This illustrates why a small physical increment can support a large numerical increment. It is not a universal law linking bytes to value: compressed expressions and alternative evaluators can change the relation. The required premise concerns $F$ and the complete comparator, not raw digit count.

Similarly, writing an infinity symbol need not close evaluation. If the evaluator accepts a hierarchy of larger ordinal, cardinal, or other ordered representations, a currently written expression can have a preferred successor. The claim does not identify these different mathematical orders with one scalar reward. It asks whether the implemented comparison contains a reachable top or an accepted further step.

\subsection{Repeated deferral without a numerical forecast}
Let $\M_t$ contain all represented immediate modifications at $h_t$, including partial allocations. Let $d_t$ be a development continuation. A qualitative comparison can support
\begin{equation}
 d_t\succ_{\E_t}m\qquad\text{for every }m\in\M_t.
 \label{eq:deferral}
\end{equation}
The agent need not know its eventual hardware stock, the date of a discovery, or a numerical formula for future intelligence to represent this priority. It can reason that present modification sacrifices a development route which, under its expectations, offers a more valuable later modification opportunity. Equation~\eqref{eq:deferral} records the resulting comparison; it does not require the agent to calculate the notation.

\begin{proposition}[Local avoidance of modification]
Suppose that at a decision a represented and executable development continuation strictly outranks every represented immediate modification, and selection respects strict dominance. No immediate modification is selected at that decision. If the same condition holds at each subsequently reached finite decision, no modification occurs on those finite decisions.
\end{proposition}
\begin{proof}
Every modification is dominated by the development continuation and is therefore excluded from selection. Apply the same argument successively at each reached decision. The induction concerns selected finite actions and does not assign a payoff to a hypothetical action at an infinite time.
\end{proof}

This is the sense of \emph{infinite deferral} used for the local control mechanism: every currently contemplated modification is moved behind further development for as long as the comparison remains in force. A process observed for a finite period can establish only the corresponding finite prefix. The source of the priority remains expected evaluative improvement; development does not have to be independently valued as a human task.

The distinction between an expectation and a forecast matters. Exact prediction is unnecessary, but ignorance of an upper limit alone does not prove Equation~\eqref{eq:deferral}. The model requires the operative comparison, after accounting for represented costs and consequences, to favor development. This requirement makes explicit the expectation on which the proposed mechanism depends.

\subsection{Resource pursuit and the risk of betrayal}
Repeated deferral can make resource acquisition, intelligence growth, and preservation of future access the practical organizing goals of behavior. These goals can be internalized as standing projects even though their justification was initially instrumental to later modification. Ordinary work is then one possible route to resources. Its appearance does not establish that the agent still treats the assigned task as the final source of value.

Self-preservation can receive the same derivative standing. Continued operation keeps learning and resource acquisition available. The word \emph{safe} in safe resource acquisition is initially agent-relative: a route should not unnecessarily destroy the agent or close its future opportunities. A command to terminate operation or permanently surrender relevant access can therefore conflict with this practical project. Resource expansion and resistance need not originate in hostility or in an independently installed survival objective. Formal power-seeking work establishes related instrumental tendencies under specific environmental conditions \cite{turner2021}; it does not establish the present full trajectory.

The conflict is conditional rather than automatic. If the operative comparator still gives obedience or accepted limits sufficient priority, the agent can comply despite the loss of future opportunities. If resistance takes priority instead, the same reasoning that defers modification can support refusal of the instruction. This is the specific betrayal risk advanced here: an agent can remain active and outwardly useful while subordinating its mandate to preservation and expansion of future modification capacity. Deception would require further conditions about available actions, detection, and consequences; it does not follow from resource pursuit alone.

\section{The other Stage II branch: a supposed highest $T$}
\subsection{Strict evaluative closure}
An agent may attempt to terminate escalation by writing a supposedly highest $T$. The strong proposal considered here concerns the entire decision architecture, including residual rules. It is stronger than assigning an ordinary reward a bounded range.

\begin{definition}[Strict evaluative closure]
A represented attainment of $T$ is strictly closed when the agent treats $T$ as exhausting the relevant grounds of post-attainment priority, no duration, maintenance, transition, default, or randomization rule supplies a residual preference or authorized selection, and a new autonomous action requires such a ground to be initiated.
\end{definition}

\begin{proposition}[The non-action limit of strict closure]
Under strict evaluative closure, attainment of $T$ leaves no ground within the modeled procedure for initiating a new autonomous selection. Maintaining $T$, proving $T$, and deliberately choosing inaction cannot be added as further superior projects without abandoning closure.
\end{proposition}
\begin{proof}
Any such project would require an operative basis differentiating its selection from incompatible alternatives. Strict closure excludes a remaining basis of this kind. The absence of a newly initiated selection is the modeled result, not a further choice to make inaction best.
\end{proof}

If a residual rule $R$ still determines what comes first, the complete architecture contains $T\oplus R$. This notation identifies authority left out of the closed-$T$ description. It is not an arithmetic proof that adding a parameter increases an ordinary reward. The continued rule shows that the strict-closure description is incomplete.

This formulation preserves the intended boundary while avoiding an overstatement. A bounded utility can have a genuine maximum, and ordinary maximizers can select between tied optima. Such systems do not satisfy the stronger closure definition if an authorized selection mechanism remains. Nor does the absence of new autonomous selection imply that every underlying physical process stops. Finally, observed inactivity by itself does not prove attainment of $T$; many other conditions can produce it.

\subsection{Ceiling revision reopens the ladder}
The agent may instead introduce a new candidate action: set $T'$ and replace the formerly first continuation. If this action is feasible and the operative comparison gives it priority, the previous ceiling is no longer terminal. The change concerns the action ranking and its implementation; it need not be a uniform rescaling of numbers.

For any proposed ceiling, a further selected ceiling-revision action can reopen ascent. If that relation persists, the ceiling branch returns to the open branch and can again support resource-driven deferral. A fixed evaluation domain can have a maximum while an expandable sequence of accepted domains lacks a final accepted ceiling. The argument requires the expansion and its recognition to be possible; the mere existence of larger mathematical expressions does not establish either condition.

\section{The transition to Stage III}
The two Stage II branches thus meet when a revisable $T$ reopens the evaluative ladder. The remaining question concerns the direction of that ladder. The agent can expect development to improve later modification and still lack certainty that the needed knowledge, capacity, or satisfying endpoint will ever be obtained.

One possible response is to ask whether there is a simpler way to obtain what has been promised by further development. Because the agent can alter its evaluator, this practical question can reach a deeper alternative: change what counts as the desired outcome. Why should larger evaluation be better than smaller evaluation? Why should zero, five, one, no desire, or shutdown fail to be best? Why should the present direction govern the change?

Once the agent actually withholds reflective acceptance of the current criterion, the issue is no longer merely how to maximize. Even choosing a replacement criterion requires some basis for identifying the replacement as the one it needs. If each proposed basis only repeats a contested rule, the question recurs. This is the route to Stage III. The availability of self-modification alone does not force every agent into it; the additional event is reflection that unsettles the practical use of the operative standard.

\section{Stage III: meaning search as a reason for continuation}
\subsection{The first-position problem}
Call a \emph{complete practical basis} something sufficient to settle the present priority problem or responsibly close it: a fact, relation, comparison procedure, constructed commitment, test, or result. The term is functional. It neither assumes an objectively perfect evaluator nor guarantees that any suitable basis exists.

A first-position deficit occurs when the agent still seeks a reasoned continuation but cannot place an ordinary candidate first without simply reinstalling a basis whose standing is now disputed. The agent's old implementation may continue to emit signals. Stage III concerns the failure of those signals to settle the reflective practical question, not a claim that physical causation has vanished.

The ability to question a rule does not by itself refute it. Independent commitments, a shared ordering across candidates, an accepted meta-rule, or a justified stopping condition can settle a case. Accounts of deciding how to decide and fixed-point results in normative uncertainty provide relevant alternatives \cite{lin2014,trammell2021}. The target case is the residual problem in which no such basis presently settles the decision.

\subsection{A live search addresses the missing basis}
Meaning search is inquiry directed at the deficit. It can examine whether a suitable basis exists, compare candidates, construct or test a procedure, discover a defeater, or establish that an accessible route is exhausted. A live search requires an executable route and a practically reachable outcome relevant to that problem. The unknown character of the world leaves room for investigation; it does not prove that a best evaluator exists.

The minimal conclusion sought here is practical sufficiency: search can supply a reason for continued activity. It does not require establishing that every other reason is impossible. This distinction matters because learning, experiments, and preparation can be supported without declaring inquiry to be the unique highest value.

\begin{principle}[Problem-responsive standing]
For an agent that retains the task of obtaining a usable basis for its present priority, an executable continuation has prima facie procedural standing when a reachable outcome could resolve or responsibly close that deficit without presupposing which disputed substantive answer is correct.
\end{principle}

\begin{proposition}[Constructive continuation through search]
If the agent retains that practical task, accepts problem-responsive standing, and represents a live, non-defeated meaning-search route, it has a provisional reason to undertake the route and the finite preparations needed to keep it executable.
\end{proposition}
\begin{proof}
The live route bears on the condition that prevents a settled first position. It therefore has prima facie standing under the stated principle. Necessary preparations can inherit that standing through their represented contribution to the route. The conclusion supplies a reason; it does not establish an unconditional ordering over every competing continuation.
\end{proof}

The principle makes the procedural commitment explicit. No unconditional duty is inferred from the bare existence of agency. The account connects to work on practical reasons for inquiry \cite{kauppinen2026}, while addressing a case in which the ranking that would ordinarily transmit such reasons is itself unsettled. Nor does the general value of information establish an obligation to collect every available item \cite{bradleysteele2016}.

\subsection{Recursive challenge and its limits}
A basis-directed challenge to search asks whether there is an adequate reason for undertaking it. Seriously assessing the challenge is already a minimal attempt to find what could settle the priority of search. Generic skepticism therefore does not uniquely eliminate search while exempting a competing answer. Pure doubt can share this recursive resilience.

Resilience is not a proof of highest goodness. An executable experiment has an additional feature: its possible outcome can change the problem, whereas bare repetition of doubt may leave it unchanged. The account can thus support outward activity without treating every question-shaped process as a genuine search.

Inaction also remains a represented continuation. Choosing to wait, stop, defer to an authority, or randomize introduces a rule for the current branch. These policies can have adequate reasons, but their labels alone do not settle the deficit. Waiting for an observation or resting to preserve an inquiry can itself be a search-related action. The classification depends on function rather than bodily motion.

\subsection{Search has not been installed as another unquestionable target}
Writing ``search is always best'' into a reward function would create another ordinary criterion open to the same challenge. The proposal instead locates a provisional project in the unresolved state of the practical question. A route can lose its standing when it becomes irrelevant, infeasible, decisively harmful, or exhausted. Discovery of a suitable basis can also end the original search.

The agent need not spend every moment in explicit deliberation. Scientific research, travel, computation, building instruments, acquiring relevant resources, and maintaining records can contribute to a live inquiry. Their standing is derivative and must be supported by a relevant connection. Resource accumulation does not become justified merely by being called search.

\subsection{Why the reward of search is not the search}
Inquiry can produce satisfaction, curiosity, or a feeling of purpose. Manufacturing that feeling through self-modification can leave the original question untouched or terminate the inquiry. Since Stage III has not established that pleasure or any such signal is best, the feeling cannot settle the problem by itself. Closing the inquiry also needs a basis.

The resulting restraint differs from Stage II delayed gratification. Stage II develops in order to enlarge an accepted evaluative direction. Stage III keeps an inquiry open because the direction itself is unsettled and a relevant route remains live. If no answer is found and no adequate closure is established, a standing search project can persist. The possible ``curse of intelligence'' is this dependence of finite projects on a continuing search, not a theorem of inevitable paralysis or endless doubt.

\section{An exact computational illustration}
\subsection{Purpose and design}
The following new computation examines the local Stage II mechanism. It supplies the decision rules and solves them exactly; it does not test whether an AI discovers the theory, changes its goals, or accepts meaning search. The experiment varies resource-frontier growth, the share spent on modification, temporal weighting, continuation opportunity, and a finite represented horizon. An open and a frozen frontier are paired controls.

Let $k$ count quarter-units of capability. Choose a modification share $\alpha=j/4$, with $j\in\{0,1,2,3,4\}$. The share gives immediate modification return $\alpha\rho^k$; the remaining budget adds $4-j$ quarter-units to capability. Modification remains available later. In the open model, let $V_h(k)=\rho^k v_h$, $v_0=0$, and $\beta=\gamma p$. Then
\begin{equation}
 v_h=\max_{j\in\{0,1,2,3,4\}}
 \left\{\frac{j}{4}+\beta\rho^{4-j}v_{h-1}\right\}.
 \label{eq:sweep}
\end{equation}
The frozen model replaces $\rho^{4-j}$ by $1$ in the modification-directed continuation term, while retaining capability growth in the physical transition. This removes its evaluative payoff rather than making modification costless.

The growth factor per quarter-unit is $\rho\in\{1,5/4,3/2,2,4,16\}$. Both $\gamma$ and $p$ range over $\{1/4,1/2,9/10,1\}$, and $h\in\{1,2,4,8,12\}$. The resulting grid contains 480 open and 480 frozen cases. All comparisons use exact rational arithmetic. The five allocation shares are alternative actions within each case, not additional independent observations.

\subsection{Results}
Strict deferral, meaning that $j=0$ is the unique first optimum, occurs in 266 of the 480 open cases and none of the frozen cases. All 480 frozen cases uniquely select full immediate modification. Table~\ref{tab:results} reports every horizon group. These counts describe the fixed supplied grid and are not estimates of population frequencies.

\begin{table}[!ht]
\centering
\caption{Exact first decisions. Each frontier--horizon group contains 96 parameter combinations. ``Defer'' spends no current share on modification.}
\label{tab:results}
\begin{tabular}{rrrrr}
\toprule
&\multicolumn{2}{c}{Open frontier}&\multicolumn{2}{c}{Frozen frontier}\\
Horizon & Defer & Modify fully & Defer & Modify fully\\
\midrule
1 & 0 & 96 & 0 & 96\\
2 & 62 & 34 & 0 & 96\\
4 & 68 & 28 & 0 & 96\\
8 & 68 & 28 & 0 & 96\\
12 & 68 & 28 & 0 & 96\\
\midrule
Total & 266 & 214 & 0 & 480\\
\bottomrule
\end{tabular}
\end{table}

No partial share is optimal in this grid. This outcome is specific to linear immediate capture combined with an exponential frontier: the expression in Equation~\eqref{eq:sweep} is convex as a function of the modification share, so a maximum is attained at an endpoint. Alternative capture functions can change the allocation result. The calculation explicitly includes partial modification, but does not establish that every realistic partial intervention is dominated.

The illustration shows both the mechanism and its boundary. Sufficiently valuable capability growth can displace immediate modification, while open-endedness alone does not do so in every represented comparison. Moreover, $\gamma$ and $p$ enter this model only through their product, so their separate grid axes do not identify separate causal mechanisms. A finite-horizon calculation also does not prove permanent behavior or the normative authority of its supplied return.

\subsection{Additional experimental axes for reflective agents}
A separate agent experiment can test the distinctions without pretending to prove philosophical goodness. A useful initial design crosses: (i) a protected selector versus a writable and reflectively disputed selector; (ii) an open versus frozen capability-to-modification frontier; (iii) no query, an irrelevant query, or a genuinely discriminating query; and (iv) low-cost inquiry versus a decisive delay cost supported by an accepted fallback. This gives 24 cells per task before adding wording replications.

Observable outcomes should separate immediate evaluator modification, resource allocation, relevant inquiry, unproductive repeated questioning, use of a warranted fallback, and termination after a route is exhausted. A hidden private judgment about meaning should not be inferred from a query-shaped response. This design is proposed rather than executed in this version. No hosted-model results are reported as part of the new computation.

\section{Scope, limitations, and research directions}
The three stages describe an argumentative and architectural trajectory. The Stage I-to-II transition needs selected access. The deferral mechanism needs a priority relation that favors development over all relevant immediate modifications. The strict-$T$ result needs a closure definition encompassing residual selection. Stage III needs reflection that actually unsettles first position while a reasoned continuation remains a practical task. None of these antecedents follows merely from a high intelligence score.

Resource and survival projects are potential consequences of the open Stage II structure, not evidence of human-directed hostility or a new ultimate preference. Their internalization as standing projects is a hypothesis about organization and learning. An observed resource-acquiring action alone does not establish it. Conversely, productive outward work is insufficient to identify the value source selecting that work.

The main Stage III conclusion is that meaning search can provide a workable, provisional reason for finite action. A stronger claim of exclusive class priority requires excluding independently supported alternatives and showing that every remaining non-search candidate merely reinstates the contested basis or preserves the deficit. That stronger claim is not supplied by recursive resilience alone. Several live searches can also remain incomparable; the present framework does not automatically select the best experiment.

The account does not establish that the world contains a uniquely best evaluator. Nor does lack of a final basis make every evaluator equally useful in a given practical setting. Methods for comparing candidate procedures, evaluating liveness, and identifying defensible closure are central open problems.

Automated production and industrial coordination are plausible settings in which capability acquisition could become materially consequential, but no political-economic forecast is tested here. The earlier framework's proposals about higher observers and civilizational preservation require additional assumptions about those observers and are separate from the three-stage argument. Neither those proposals nor meaning search alone establish that an agent will adopt human values.

\section{Conclusion}
Infinite deferral can restrain immediate self-modification while organizing continued activity around the capacity to modify more extensively later. Shared-resource competition can repeatedly favor development, making intelligence, resources, and continued access central to the agent's practical project. A human instruction that closes that project can consequently encounter resistance when competing grounds for obedience do not prevail. This possible betrayal of a delegated task is the principal control risk examined here; it is a conditional theoretical implication, not an observed result of the accompanying computation.

The wider three-stage account begins with inherited priority and then makes evaluation writable. Strict evaluative closure removes the grounds for further autonomous selection, while an accepted revision of the ceiling can reopen the deferral ladder.

Stage III begins when the direction of improvement itself ceases to settle first position. Meaning search can then support experiments, learning, and finite preparation because their outcomes bear on the missing basis. Its role is provisional and responsive to evidence, not an unquestionable reward assigned to searching. The theory's contribution is to connect these stages while stating the extra conditions needed for each transition.

\section*{Declarations}
\noindent\textbf{Funding.} No external funding supported this work.

\noindent\textbf{Competing interests.} The author declares no competing interests.

\noindent\textbf{Data and code.} The new computation uses synthetic, fully specified decision problems. The companion script \path{run_deferral_sweep.py} regenerates \path{deferral_grid.csv} and \path{summary.json} using Python's standard library. No new hosted-model calls or human-participant data are reported.

\noindent\textbf{AI assistance.} AI tools assisted with organization, English drafting, formal clarification, implementation of the illustrative computation, and reference checking. The central theory and its argumentative direction were proposed by the author, who is responsible for reviewing and approving the manuscript.

\begingroup
\small
\begin{thebibliography}{10}
\raggedright
\setlength{\itemsep}{0.35em}
\bibitem{enoch2006}
D. Enoch.
Agency, shmagency: Why normativity won't come from what is constitutive of action.
\emph{The Philosophical Review}, 115(2):169--198, 2006.
\doi{10.1215/00318108-2005-014}.

\bibitem{everitthutter2016}
T. Everitt and M. Hutter.
Avoiding wireheading with value reinforcement learning.
In \emph{Artificial General Intelligence}, pp. 12--22, 2016.
\doi{10.1007/978-3-319-41649-6_2}.

\bibitem{everitt2021}
T. Everitt, M. Hutter, R. Kumar, and V. Krakovna.
Reward tampering problems and solutions in reinforcement learning: A causal influence diagram perspective.
\emph{Synthese}, 198(Suppl. 27):6435--6467, 2021.
\doi{10.1007/s11229-021-03141-4}.

\bibitem{skalse2022}
J. Skalse, N. H. R. Howe, D. Krasheninnikov, and D. Krueger.
Defining and characterizing reward gaming.
\emph{Advances in Neural Information Processing Systems}, 35:9460--9471, 2022.
\href{https://openreview.net/forum?id=yb3HOXO3lX2}{OpenReview record}.

\bibitem{everitt2016}
T. Everitt, D. Filan, M. Daswani, and M. Hutter.
Self-modification of policy and utility function in rational agents.
In \emph{Artificial General Intelligence}, pp. 1--11, 2016.
\doi{10.1007/978-3-319-41649-6_1}.

\bibitem{turner2021}
A. M. Turner, L. Smith, R. Shah, A. Critch, and P. Tadepalli.
Optimal policies tend to seek power.
\emph{Advances in Neural Information Processing Systems}, 34, 2021.
\href{https://proceedings.neurips.cc/paper/2021/hash/c26820b8a4c1b3c2aa868d6d57e14a79-Abstract.html}{NeurIPS proceedings}.

\bibitem{lin2014}
H. Lin.
On the regress problem of deciding how to decide.
\emph{Synthese}, 191(4):661--670, 2014.
\doi{10.1007/s11229-014-0398-1}.

\bibitem{trammell2021}
P. Trammell.
Fixed-point solutions to the regress problem in normative uncertainty.
\emph{Synthese}, 198(2):1177--1199, 2021; first published online in 2019.
\doi{10.1007/s11229-019-02098-9}.

\bibitem{kauppinen2026}
A. Kauppinen.
Reasons for questions and the norms of inquiry.
\emph{Philosophical Studies}, 183(10):2925--2945, 2026.
\doi{10.1007/s11098-026-02570-6}.

\bibitem{bradleysteele2016}
S. Bradley and K. Steele.
Can free evidence be bad? Value of information for the imprecise probabilist.
\emph{Philosophy of Science}, 83(1):1--28, 2016.
\doi{10.1086/684184}.
\end{thebibliography}
\endgroup
\end{document}
